\section{A full cross-flight root sum from curve growth}
\label{sec:ag-growth}
We spell out the remaining dynamic input, including its curve class.
This prevents replacement of a global pressure premise by an unstated
regularity premise. The exact geometry of the preceding section is
unconditional within the retained physical finite-horizon model.

\begin{inputstatement}[Homogeneous curve growth, in its required class]
\label{input:ag-growth}
Choose the common cone upper bound
$U=\kappa_++\tau_-^{-1}$ and a regular homogeneous unstable
curve class containing all the graphs of
Theorem~\ref{thm:ag-family}. In particular its curvature allowance
must accommodate $K_0$; being merely $C^2$ does not establish this
membership. For some common $k_0\ge2$, sufficiently small
$\ell_*>0$, and $C_g\ge1$, every curve $C$ of that class and
length at most $\ell_*$ has a source-labeled partition of its
regular forward image into homogeneous short curves
$\mathcal G^+_l(C)$ satisfying
\begin{equation}\label{eq:ag-growth}
 \sum_{V\in\mathcal G^+_l(C)}
       \sup_V J_V(T_R^{-l})\le C_g,
             \qquad R\in[R_-,R_+],\quad l\ge0.
\end{equation}
The branches cover arclength-almost every original point, the
inverse branch in each summand is the actual physical one, and
all resulting graphs have slopes in $[\kappa_-,U]$.
The Jacobian is for Euclidean arclength in $(s,\varphi)$.
\end{inputstatement}

This is the forward, time-reversed version of the curve input in
Section~\ref{sec:cf-transfer}, with the source-class membership
now made explicit. The relevant published curve-sum statement is
\cite[Lemma 3.2(b)]{DZUniform}; it is based on the homogeneous curve,
expansion and distortion properties (H2)--(H4), not on a sum of
physical cylinder suprema. The stated class and its family-uniform
constants remain explicit external inputs here. No independent
reconstruction of the full geometric-pressure theorem, nor of the
older singular pairing and spectral interfaces, is claimed.

Shrinking $\ell_*$ if necessary, require
\begin{equation}\label{eq:ag-length-choice}
 \ell_*\le\frac{\pi}{2D_0},\qquad
 D_0=\kappa_++U+(1-R_+)^{-1}.
\end{equation}
The admissible growth partition uses this common choice; any change
of its subdivision scale must be included in its constant $C_g$.

\begin{lemma}[Counted clearance on a short unstable graph]
\label{lem:ag-band}
For a Borel set $I\subset(0,\delta_*)$, let $N_I(x)$ count all
unvisited disks on the next actual flight of $x$ whose clearance
belongs to $I$. For every graph $V$ in
Input~\ref{input:ag-growth} of length at most $\ell_*$,
\begin{equation}\label{eq:ag-band}
 \int_V N_I\,d\ell\le C_t|I|,
 \qquad C_t=\frac{2N_\circ\sqrt{1+U^2}}{\kappa_-a_*}.
\end{equation}
This is a bound on the mark \emph{count}, including both impact
signs, not just on its nonzero indicator.
\end{lemma}
\begin{proof}
Use a lift with the incident disk centered at zero. Every possibly
marked center belongs to $\mathcal Z_\circ$. For a fixed such $z$,
let $\vartheta(s)$ be the outgoing direction, and put
$b_z=(z-q)\cdot v_\vartheta^\perp$,
$t_z=(z-q)\cdot v_\vartheta$. Direct differentiation gives
\[
 b_z'=-\cos\varphi-t_z\vartheta',\qquad
                 \vartheta'=1/R+\varphi'\ge\kappa_-.
\]
The angle $\eta=\arg(z-q)-\vartheta$ is strictly decreasing:
$\eta'\le(1-R_+)^{-1}-\kappa_-<0$. Its total variation is at
most $D_0\ell_*\le\pi/2$. Thus $t_z>0$, equivalently
$\cos\eta>0$, holds on at most one interval of the graph.
On this interval $b_z$ is strictly decreasing. On the physical
mark subset the interior-foot estimate gives $t_z\ge a_*$;
therefore $|b_z'|\ge\kappa_-a_*$. For each sign, change variables
in this one-to-one map on the marked subset. The integral of its
indicator is at most
$\sqrt{1+U^2}|I|/(\kappa_-a_*)$. Dropping a first-hit or
unvisited mask occurs only in this positive bound. Sum the two
signs and the $N_\circ$ candidates. Approximation by open sets,
or the one-dimensional area formula for a monotone smooth map,
gives the claim for every Borel $I$. No number of components of
$I$ or of an exact-event selector occurs.
\end{proof}

\begin{theorem}[A fixed-first-clearance kernel bound]
\label{thm:ag-kernel}
Assume Input~\ref{input:ag-growth}. For every $R\in[R_-,R_+]$,
every fixed $0<u_1\le\delta_*$, every integer $h\ge1$, and every
Borel $I\subset(0,\delta_*)$,
\begin{equation}\label{eq:ag-kernel}
 \sum_{\sigma=\pm1}\int_0^{2\pi}
 N_I\bigl(T_R^hX_R(\vartheta,\sigma,u_1)\bigr)\,d\vartheta
       \le H_*|I|,\qquad H_*=A_*C_gC_t.
\end{equation}
The integrand is taken on its actual regular domain. Its exceptional
set has zero angular measure under the curve input. Arbitrary Borel
weights between zero and one, including the original section, exact
record and physical roof selectors, may be retained inside the sum
without increasing the bound. In particular the second-clearance
measure has a density bounded by $H_*$, uniformly in the first
clearance and in the complete collision separation.
\end{theorem}
\begin{proof}
The first actual collision has already been taken in
$Y=T_RX$. Apply Theorem~\ref{thm:ag-family} to its source.
Exactly $l=h-1$ further iterates remain. For each labeled curve
$C$, use the physical image partition in
\eqref{eq:ag-growth} and one-dimensional change of variables:
\begin{align*}
 \int_C \rho_C(y)N_I(T_R^ly)\,d\ell_C(y)
 &=\sum_{V\in\mathcal G_l^+(C)}\int_V
    (\rho_C\circ T_R^{-l})N_I\,J_VT_R^{-l}\,d\ell_V\!\!\\
 &\le \|\rho_C\|_\infty C_g C_t |I|.
\end{align*}
Use Lemma~\ref{lem:ag-band} separately on each $V$, and then sum
its weighted inverse Jacobian, not its number of components.
Summing $C$ uses \eqref{eq:ag-family} and proves the bound. For
$h=1$ the middle map is the identity and $l=0$; no collision is
missing or added. First truncate the nonnegative sums and then
use monotone convergence. Carry a bounded Borel selector through
each exact substitution before dropping it in the inequality.
The measure domination for all $I$ implies absolute continuity
and the displayed density bound by the Radon--Nikodym theorem.
It concerns clearance, not the roof variable $L_{m,R}$.
\end{proof}

\begin{theorem}[One-copy marked root pressure without physical-word pressure]
\label{thm:ag-roots}
Under precisely the preceding curve-growth input, let
$\mathcal A_{h,R}(u_1,u_2)$ be the set of actual tangentially
augmented itineraries in Theorem~\ref{thm:xf-coarea}. Then
\begin{equation}\label{eq:ag-roots}
 \sup_{R\in[R_-,R_+],\ h\ge1}\ 
 \esssup_{(u_1,u_2)\in(0,\delta_*)^2}
       \sum_{a\in\mathcal A_{h,R}(u_1,u_2)}B_h(x_a)^{-1}
                  \le H_*.
\end{equation}
The complete positive physical sum precedes the essential supremum.
No premise on $\mathcal P_l(R)$, no product of independently chosen
flights and no growth-rate bound on the number of words is used.
Consequently
\begin{equation}\label{eq:ag-pair}
 \int N_{\delta_1}(N_{\delta_2}\circ T_R^h)\,d\nu
       \le C_{\rm ag}\delta_1\delta_2,
       \qquad C_{\rm ag}=H_*/(4\pi R_-).
\end{equation}
\end{theorem}
\begin{proof}
Use the exact one-mark flux, followed by the actual
cross-flight coarea identity of Theorem~\ref{thm:xf-coarea}.
Global one-copy multiplicity is supplied by
Theorem~\ref{thm:xf-unique}, not by local invertibility. For an
arbitrary Borel set $I_1\subset(0,\delta_*)$ and $I_2$ in the
same interval, the resulting joint clearance measure is both
\[
 \int_{I_1}\sum_{\sigma}\int_0^{2\pi}
 N_{I_2}(T_R^hX_R(\vartheta,\sigma,u_1))\,d\vartheta\,du_1
 \quad\hbox{and}\quad
 \int_{I_1\times I_2}\sum_{a\in\mathcal A_{h,R}}B_h(x_a)^{-1}
                                                    \,du_1du_2.
\]
The common collision prefactor $1/(4\pi R)$ is omitted on
\emph{both} sides of this equality. Theorem~\ref{thm:ag-kernel}
bounds this measure by $H_*du_1du_2$, first on rectangles and
then on all Borel sets. Uniqueness of its density proves
\eqref{eq:ag-roots}. Reintroducing $1/(4\pi R)$ and integrating
over the two clearance intervals proves \eqref{eq:ag-pair}.
For each fixed radius the countable separations admit a common
clearance-null exceptional set. A null set common to an uncountable
radius family is neither needed nor asserted.
\end{proof}

\subsection{The inherited complete-record conclusion under the curve input}
The conclusion below was already obtained in
Theorem~\ref{thm:xf-full-count} with the physical pressure input.
The new contribution is its proof from the directly lifted angular
source and curve growth, and the stronger fixed-first-clearance
kernel bound, not a second claim of novelty for the integer identity.

\begin{corollary}[Complete chronological multiplicity with the revised dependency]
\label{cor:ag-record}
Put $C_A=\max\{2C_{\rm pair},C_{\rm ag}\}$, where $C_{\rm pair}$
is the inherited same-flight constant. For any finite deterministic
sequence $0\le\delta_j\le\delta_*$, put
$\mathscr N=\sum_{j=0}^{m-1}N_{\delta_j}(T_R^jx)$ and
$S=\sum_j\delta_j$. Under Input~\ref{input:ag-growth},
\begin{equation}\label{eq:ag-factorial}
 \int\mathscr N\,d\nu=S/R,\qquad
 \int\mathscr N(\mathscr N-1)\,d\nu\le C_A S^2.
\end{equation}
The first equality requires no curve-growth assumption. For the
original exact event $\Pi$ and any Borel $0\le\psi\le1$,
\begin{equation}\label{eq:ag-exact-source}
 0\le\int\one_\Pi\psi(L_{m,R})
          [\mathscr N-\one_{\{\mathscr N\ge1\}}]d\nu_R^*
              \le\frac{C_A S^2}{2c_*}.
\end{equation}
In particular, for $\delta_j=2\varepsilon r^{\min(j,m-1-j)}$,
$r=q_0^{1/4}$ and $2\varepsilon\le\delta_*$, this is at most
$8C_A\varepsilon^2/[c_*(1-r)^2]$ for every $m$.
\end{corollary}
\begin{proof}
Invariance of $\nu$ and the exact one-mark mean give the first
identity. Expand the factorial exactly as
\[
 \mathscr N(\mathscr N-1)=\sum_jN_j(N_j-1)
                         +2\sum_{j<k}N_jN_k.
\]
The same-flight theorem bounds the first sum by
$2C_{\rm pair}\sum_j\delta_j^2$ after integration. Invariance
and \eqref{eq:ag-pair} bound the second by
$2C_{\rm ag}\sum_{j<k}\delta_j\delta_k$. Their sum is at most
$C_A S^2$. No independence was used. The deterministic inequality
$0\le n-\one_{\{n\ge1\}}\le n(n-1)/2$ for integers $n\ge0$
and $\Pi\subset Y_R^*$ give \eqref{eq:ag-exact-source}; only here
is the source dominated by $\nu/c_*$. There are at most two
indices at each depth, so $S\le4\varepsilon/(1-r)$, proving the
last bound. Repeated sightings of a disk on distinct flights are
included in the cross terms.
\end{proof}

\subsection{The original roof obstruction has not been replaced}
The new norm in \eqref{eq:ag-family} describes a singular
one-dimensional angular source in collision space. The density
in \eqref{eq:ag-kernel} is in \emph{second clearance}. It is not
a density in the original roof, and the complete root bound in
\eqref{eq:ag-roots} has no hidden $m^{-2}$ factor. The unconditioned
first-order mass $S/R$ remains positive. An exact-label selection
can change its leading coefficient; it is not replaced by a
probability conditioned on a mark.

To reach the raw target, the original ordered shortage limit
\eqref{eq:v112-remaining} must still be proved. In particular the
first-order marked source must be estimated after the complete
original roof pushforward and physical word sum, including fold,
critical, higher-degeneracy, one-sided and pairing-boundary terms,
with the required $m^2$ normalization. The single-mark fixed-clearance
roof bridge in \eqref{eq:cf-bridge-input} remains only a sufficient
unproved condition; it is not a premise added to our main problem.
The local fold calculations show why it is safer to integrate
clearance before taking a roof supremum. Low-incidence shallow
source and nongrazing capacity likewise remain in the original
shortage. At fixed outer $B$, multiplying
\eqref{eq:ag-exact-source} by $m^2$ does not prove decay. No
collision-dependent roof window, altered section, exceptional-set
subtraction or changed order of limits is introduced here.

